% GENERATED FILE. Edit canonical lessons, then run npm run generate:books.
% canonical-source-hash: fnv1a64:8b868b801716ed4a
% canonical-lessons: HI-C382-jaldbazi, HI-C382-laparvah, HI-C382-chalakna, HI-C382-masalna, HI-C382-nazuk, HI-R382-review-375-382

\chapter{Haste and Care}
\label{ch:hi-haste-and-care}
\hlchaptermodality{382}

\begin{chapteropening}
\textbf{By the end of this chapter:} \emph{I can warn someone what will happen if they rush.}

Complete the last lesson of chapter 382: I can warn someone what will happen if they rush.
\end{chapteropening}

\section[\texorpdfstring{\hi{जल्दबाज़ी} (jaldbāzī)}{जल्दबाज़ी (jaldbāzī)}]{\hi{जल्दबाज़ी} (jaldbāzī) — haste, rushing}

\label{lesson:HI-C382-jaldbazi}



\begin{quote}
Before the new one: say the Hindi for to waste, then the Hindi for decent.
\end{quote}

\subsection*{You'll want to know: \hi{जल्दबाज़ी}}
\textbf{\hi{जल्दबाज़ी}} — \emph{jaldbāzī} — \textquotedblleft{}haste, rushing\textquotedblright{}.

\begin{quote}
\hi{जल्दबाज़ी} \hi{मत} \hi{करो}, \hi{वरना} \hi{ग़लती} \hi{होगी।}
\end{quote}

\begin{quote}
\emph{jaldbāzī mat karo, varnā ġaltī hogī.}
\end{quote}

\begin{quote}
Don't rush, or you'll make a mistake.
\end{quote}

Feminine.

\begin{cousinweb}
\textbf{\hi{जल्दबाज़ी}} (\emph{jaldbāzī}) is usually traced to Persian \emph{jald}, \textquotedblleft{}quick\textquotedblright{}, and \emph{bāzī}, \textquotedblleft{}play, a game\textquotedblright{}: playing at speed.
\end{cousinweb}

\begin{grammarlens}[title={What you've built}]
One more word for saying what happens if.
\end{grammarlens}

\subsection*{Your turn}
\begin{itemize}
\raggedright
  \item \emph{Say it:} \emph{jaldbāzī}
  \item \emph{Say it:} \emph{jaldbāzī}, once more
  \item \emph{Say it:} \emph{jaldbāzī mat karo, varnā ġaltī hogī.} — Don't rush, or you'll make a mistake.
  \item \emph{Recall:} say the Hindi for to waste, then the Hindi for decent, then say \emph{jaldbāzī} again
\end{itemize}

\begin{tcolorbox}[breakable,colback=teal!4,colframe=teal!35!black,title={Before you move on}]
What does \hi{जल्दबाज़ी} mean? (\textquotedblleft{}Haste, rushing\textquotedblright{}.) Say it once more.
\end{tcolorbox}
\section[\texorpdfstring{\hi{लापरवाह} (lāparvāh)}{लापरवाह (lāparvāh)}]{\hi{लापरवाह} (lāparvāh) — careless}

\label{lesson:HI-C382-laparvah}



\begin{quote}
Before the new one: say the Hindi for decent, then the Hindi for haste.
\end{quote}

\subsection*{You'll want to know: \hi{लापरवाह}}
\textbf{\hi{लापरवाह}} — \emph{lāparvāh} — \textquotedblleft{}careless\textquotedblright{}.

\begin{quote}
\hi{अगर} \hi{ड्राइवर} \hi{लापरवाह} \hi{होगा} \hi{तो} \hi{दुर्घटना} \hi{हो} \hi{सकती} \hi{है।}
\end{quote}

\begin{quote}
\emph{agar ḍrāivar lāparvāh hogā to durghaṭnā ho saktī hai.}
\end{quote}

\begin{quote}
If the driver is careless, there can be an accident.
\end{quote}

\begin{cousinweb}
\textbf{\hi{लापरवाह}} (\emph{lāparvāh}) is usually traced to Persian \emph{lāparvā}, \textquotedblleft{}heedless\textquotedblright{}.
\end{cousinweb}

\begin{grammarlens}[title={What you've built}]
One more word for saying what happens if.
\end{grammarlens}

\subsection*{Your turn}
\begin{itemize}
\raggedright
  \item \emph{Say it:} \emph{lāparvāh}
  \item \emph{Say it:} \emph{lāparvāh}, once more
  \item \emph{Say it:} \emph{agar ḍrāivar lāparvāh hogā to durghaṭnā ho saktī hai.} — If the driver is careless, there can be an accident.
  \item \emph{Recall:} say the Hindi for decent, then the Hindi for haste, then say \emph{lāparvāh} again
\end{itemize}

\begin{tcolorbox}[breakable,colback=teal!4,colframe=teal!35!black,title={Before you move on}]
What does \hi{लापरवाह} mean? (\textquotedblleft{}Careless\textquotedblright{}.) Say it once more.
\end{tcolorbox}
\section[\texorpdfstring{\hi{छलकना} (chalaknā)}{छलकना (chalaknā)}]{\hi{छलकना} (chalaknā) — to spill over, to slop}

\label{lesson:HI-C382-chalakna}



\begin{quote}
Before the new one: say the Hindi for haste, then the Hindi for careless.
\end{quote}

\subsection*{You'll want to know: \hi{छलकना}}
\textbf{\hi{छलकना}} — \emph{chalaknā} — \textquotedblleft{}to spill over, to slop\textquotedblright{}.

\begin{quote}
\hi{गिलास} \hi{धीरे} \hi{से} \hi{उठाओ}, \hi{वरना} \hi{पानी} \hi{छलक} \hi{जाएगा।}
\end{quote}

\begin{quote}
\emph{gilās dhīre se uṭhāo, varnā pānī chalak jāegā.}
\end{quote}

\begin{quote}
Lift the glass slowly, or the water will spill.
\end{quote}

Past (it …): \hi{वह} \hi{दूध} \hi{छलक} \hi{गया}, \emph{vah dūdh chalak gayā}.

\begin{cousinweb}
\textbf{\hi{छलकना}} (\emph{chalaknā}) is inherited from \emph{*chala}, \textquotedblleft{}a splash\textquotedblright{}, a form CDIAL reconstructs from its descendants (CDIAL 5002).
\end{cousinweb}

\begin{grammarlens}[title={What you've built}]
One more word for saying what happens if.
\end{grammarlens}

\subsection*{Your turn}
\begin{itemize}
\raggedright
  \item \emph{Say it:} \emph{chalaknā}
  \item \emph{Say it:} \emph{chalaknā}, once more
  \item \emph{Say it:} \emph{gilās dhīre se uṭhāo, varnā pānī chalak jāegā.} — Lift the glass slowly, or the water will spill.
  \item \emph{Recall:} say the Hindi for haste, then the Hindi for careless, then say \emph{chalaknā} again
\end{itemize}

\begin{tcolorbox}[breakable,colback=teal!4,colframe=teal!35!black,title={Before you move on}]
What does \hi{छलकना} mean? (\textquotedblleft{}To spill over, to slop\textquotedblright{}.) Say it once more.
\end{tcolorbox}
\section[\texorpdfstring{\hi{मसलना} (masalnā)}{मसलना (masalnā)}]{\hi{मसलना} (masalnā) — to rub hard, to crush}

\label{lesson:HI-C382-masalna}



\begin{quote}
Before the new one: say the Hindi for careless, then the Hindi for to spill over.
\end{quote}

\subsection*{You'll want to know: \hi{मसलना}}
\textbf{\hi{मसलना}} — \emph{masalnā} — \textquotedblleft{}to rub hard, to crush\textquotedblright{}.

\begin{quote}
\hi{आँखें} \hi{मत} \hi{मसलो}, \hi{वरना} \hi{वे} \hi{और} \hi{लाल} \hi{हो} \hi{जाएँगी।}
\end{quote}

\begin{quote}
\emph{āṅkheṁ mat maslo, varnā ve aur lāl ho jāeṅgī.}
\end{quote}

\begin{quote}
Don't rub your eyes, or they'll get even redder.
\end{quote}

Past, with \hi{ने} (he or she …): \hi{उसने} \hi{फूल} \hi{मसला}, \emph{usne phūl maslā}.

\begin{cousinweb}
\textbf{\hi{मसलना}} (\emph{masalnā}) is a tadbhava, inherited from Sanskrit \emph{maṣati}, \textquotedblleft{}rubs, crushes\textquotedblright{} (CDIAL 9919).
\end{cousinweb}

\begin{grammarlens}[title={What you've built}]
One more word for saying what happens if.
\end{grammarlens}

\subsection*{Your turn}
\begin{itemize}
\raggedright
  \item \emph{Say it:} \emph{masalnā}
  \item \emph{Say it:} \emph{masalnā}, once more
  \item \emph{Say it:} \emph{āṅkheṁ mat maslo, varnā ve aur lāl ho jāeṅgī.} — Don't rub your eyes, or they'll get even redder.
  \item \emph{Recall:} say the Hindi for careless, then the Hindi for to spill over, then say \emph{masalnā} again
\end{itemize}

\begin{tcolorbox}[breakable,colback=teal!4,colframe=teal!35!black,title={Before you move on}]
What does \hi{मसलना} mean? (\textquotedblleft{}To rub hard, to crush\textquotedblright{}.) Say it once more.
\end{tcolorbox}
\section[\texorpdfstring{\hi{नाज़ुक} (nāzuk)}{नाज़ुक (nāzuk)}]{\hi{नाज़ुक} (nāzuk) — delicate, fragile}

\label{lesson:HI-C382-nazuk}



\begin{quote}
Before the new one: say the Hindi for to spill over, then the Hindi for to rub hard.
\end{quote}

\subsection*{You'll want to know: \hi{नाज़ुक}}
\textbf{\hi{नाज़ुक}} — \emph{nāzuk} — \textquotedblleft{}delicate, fragile\textquotedblright{}.

\begin{quote}
\hi{यह} \hi{शीशा} \hi{नाज़ुक} \hi{है}, \hi{अगर} \hi{गिरेगा} \hi{तो} \hi{टूट} \hi{जाएगा।}
\end{quote}

\begin{quote}
\emph{yah śīśā nāzuk hai, agar giregā to ṭūṭ jāegā.}
\end{quote}

\begin{quote}
This glass is fragile; if it falls, it will break.
\end{quote}

\begin{cousinweb}
\textbf{\hi{नाज़ुक}} (\emph{nāzuk}) is usually traced to Persian \emph{nāzuk}, \textquotedblleft{}delicate, slender\textquotedblright{}.
\end{cousinweb}

\begin{grammarlens}[title={What you've built}]
One more word for saying what happens if.
\end{grammarlens}

\subsection*{Your turn}
\begin{itemize}
\raggedright
  \item \emph{Say it:} \emph{nāzuk}
  \item \emph{Say it:} \emph{nāzuk}, once more
  \item \emph{Say it:} \emph{yah śīśā nāzuk hai, agar giregā to ṭūṭ jāegā.} — This glass is fragile; if it falls, it will break.
  \item \emph{Recall:} say the Hindi for to spill over, then the Hindi for to rub hard, then say \emph{nāzuk} again
\end{itemize}

\begin{tcolorbox}[breakable,colback=teal!4,colframe=teal!35!black,title={Before you move on}]
What does \hi{नाज़ुक} mean? (\textquotedblleft{}Delicate, fragile\textquotedblright{}.) Say it once more.
\end{tcolorbox}
\section[(dialogue)]{Looking back: from Hopes and Fears to Haste and Care}

\label{lesson:HI-R382-review-375-382}



\begin{quote}
Say the Hindi for an expectation and an apprehension.
\end{quote}

\subsection*{Your turn}
Say these aloud:

\begin{itemize}
\raggedright
  \item an expectation, an apprehension, an indication, a hint, to charm
  \item evident, definite, without fail, bound to happen, to emerge
  \item a situation, a state (of health, repair), a need, a hardship, to trap
  \item essential, compulsory, solely, except, to forbid
  \item if (formal), otherwise (formal), mere, beware, to slam
  \item a hassle, an obstacle, a solution, to get stuck, to pull out of trouble
  \item if only, a pity, to torment, to waste, decent
  \item haste, careless, to spill over, to rub hard, fragile
\end{itemize}

\begin{tcolorbox}[breakable,colback=teal!4,colframe=teal!35!black,title={Before you move on}]
An expectation? (\textbf{\hi{आशा}}, \emph{āśā}.) An apprehension? (\textbf{\hi{आशंका}}, \emph{āśaṅkā}.) An indication? (\textbf{\hi{संकेत}}, \emph{saṅket}.) A hint? (\textbf{\hi{इशारा}}, \emph{iśārā}.) To charm? (\textbf{\hi{रिझाना}}, \emph{rijhānā}.) Evident? (\textbf{\hi{ज़ाहिर}}, \emph{zāhir}.) Definite? (\textbf{\hi{निश्चित}}, \emph{niścit}.) Without fail? (\textbf{\hi{अवश्य}}, \emph{avaśya}.) Bound to happen? (\textbf{\hi{लाज़िमी}}, \emph{lāzimī}.) To emerge? (\textbf{\hi{उभरना}}, \emph{ubharnā}.) A situation? (\textbf{\hi{स्थिति}}, \emph{sthiti}.) A state (of health, repair)? (\textbf{\hi{हालत}}, \emph{hālat}.) A need? (\textbf{\hi{ज़रूरत}}, \emph{zarūrat}.) A hardship? (\textbf{\hi{मुसीबत}}, \emph{musībat}.) To trap? (\textbf{\hi{फँसाना}}, \emph{phãsānā}.) Essential? (\textbf{\hi{आवश्यक}}, \emph{āvaśyak}.) Compulsory? (\textbf{\hi{अनिवार्य}}, \emph{anivārya}.) Solely? (\textbf{\hi{केवल}}, \emph{keval}.) Except? (\textbf{\hi{सिवाय}}, \emph{sivāy}.) To forbid? (\textbf{\hi{मना} \hi{करना}}, \emph{manā karnā}.) If (formal)? (\textbf{\hi{यदि}}, \emph{yadi}.) Otherwise (formal)? (\textbf{\hi{अन्यथा}}, \emph{anyathā}.) Mere? (\textbf{\hi{मात्र}}, \emph{mātra}.) Beware? (\textbf{\hi{ख़बरदार}}, \emph{khabardār}.) To slam? (\textbf{\hi{पटकना}}, \emph{paṭaknā}.) A hassle? (\textbf{\hi{दिक़्क़त}}, \emph{diqqat}.) An obstacle? (\textbf{\hi{बाधा}}, \emph{bādhā}.) A solution? (\textbf{\hi{हल}}, \emph{hal}.) To get stuck? (\textbf{\hi{अटकना}}, \emph{aṭaknā}.) To pull out of trouble? (\textbf{\hi{उबारना}}, \emph{ubārnā}.) If only? (\textbf{\hi{काश}}, \emph{kāś}.) A pity? (\textbf{\hi{अफ़सोस}}, \emph{afsos}.) To torment? (\textbf{\hi{सताना}}, \emph{satānā}.) To waste? (\textbf{\hi{गँवाना}}, \emph{gãvānā}.) Decent? (\textbf{\hi{भला}}, \emph{bhalā}.) Haste? (\textbf{\hi{जल्दबाज़ी}}, \emph{jaldbāzī}.) Careless? (\textbf{\hi{लापरवाह}}, \emph{lāparvāh}.) To spill over? (\textbf{\hi{छलकना}}, \emph{chalaknā}.) To rub hard? (\textbf{\hi{मसलना}}, \emph{masalnā}.) Fragile? (\textbf{\hi{नाज़ुक}}, \emph{nāzuk}.)
\end{tcolorbox}
